\begin{derivation}[从\eqnrefs{eq:pmns-definition-dp11}到\eqnrefs{eq:neutrino-oscillation-probability-dp14}]
Takagi 分解后，味本征场与质量本征场的左手分量满足

\begin{equation*}
\nu_{\alpha L}=\sum_{i=1}^3U_{\alpha i}\nu_{iL},
\qquad U\equiv U_{\rm PMNS}.
\end{equation*}
因此在产生点制备的味态为
\begin{equation*}
|\nu_\alpha\rangle=\sum_iU_{\alpha i}^*|\nu_i\rangle.
\end{equation*}

各质量分量具有共同主导能量 $E$ 且满足 $m_i^2c^4/E^2\ll1$ 时，精确色散关系
\begin{equation*}
E_i^2=p_i^2c^2+m_i^2c^4
\end{equation*}
给出
\begin{align*}
p_ic
&=E\sqrt{1-\frac{m_i^2c^4}{E^2}}\\
&=E\left[1-\frac{m_i^2c^4}{2E^2}
+\mathcal O\!\left(\frac{m_i^4c^8}{E^4}\right)\right],
\end{align*}
所以
\begin{equation*}
p_i=\frac Ec-\frac{m_i^2c^3}{2E}
+\mathcal O\!\left(\frac{m_i^4c^7}{E^3}\right).
\end{equation*}
传播距离 $L$ 的平面波相位为
\begin{equation*}
\exp\!\left[-\frac{\ii}{\hbar}(E_it-p_iL)\right].
\end{equation*}
在各波包仍相干重叠时取 $t=L/c$ 到主导阶，并去掉所有质量分量共有的整体相位，留下
\begin{equation*}
\exp\!\left[-\ii\frac{m_i^2c^3L}{2E\hbar}\right].
\end{equation*}
于是
\begin{equation*}
\mathcal A_{\alpha\to\beta}(L)
=\sum_iU_{\beta i}U_{\alpha i}^*
\exp\!\left[-\ii\frac{m_i^2c^3L}{2E\hbar}\right].
\end{equation*}

取模平方并把 $i=j$ 与 $i\neq j$ 分开：
\begin{align*}
P_{\alpha\to\beta}
={}&\sum_i|U_{\beta i}|^2|U_{\alpha i}|^2\\
&+2\Re\sum_{i<j}
U_{\beta i}U_{\alpha i}^*U_{\beta j}^*U_{\alpha j}
\exp\!\left[-\ii\frac{\Delta m_{ij}^2c^3L}{2E\hbar}\right].
\end{align*}
利用幺正性 $\sum_iU_{\beta i}U_{\alpha i}^*=\delta_{\alpha\beta}$，并使用 $K_{ij}^{\alpha\beta}$ 与 $\Delta_{ij}$ 的定义，将指数的实部和虚部分开，即得到\eqnrefs{eq:neutrino-oscillation-probability-dp14}。

\end{derivation}

\begin{derivation}[从\eqnrefs{eq:smeft-geometric-normalization-dp13}到\eqnrefs{eq:heavy-scalar-ch-loop-dp32}]
取几何量纲为 $\mathrm{m^{-1}}$ 的重实标量 $S$，并令

\begin{equation*}
\widehat\Lag_S
=\frac12(\partial_\mu S)(\partial^\mu S)
-\frac12\Lambda_*^2S^2
-\frac\kappa2S^2(\widehat\Phi^\dagger\widehat\Phi),
\end{equation*}
其中 $\kappa$ 无量纲。经典分支 $S=0$ 不产生 $\mathcal O_H$，所以最低非零项来自高斯泛函积分的一圈行列式。转到欧几里得形式并令 $\widehat H=\widehat\Phi^\dagger\widehat\Phi$：
\begin{equation*}
\frac{\Delta\Gamma_E^{(1)}}{\hbar}
=\frac12\Tr\ln[-\partial_E^2+\Lambda_*^2+\kappa\widehat H].
\end{equation*}
在缓慢变化背景下忽略 $\partial\widehat H$，把场依赖部分写成
\begin{equation*}
\frac12\int\frac{\dd^4 k_E}{(2\pi)^4}
\ln\!\left(1+\frac{\kappa\widehat H}{k_E^2+\Lambda_*^2}\right).
\end{equation*}
展开 $\ln(1+u)=u-u^2/2+u^3/3+\cdots$。前两项分别重整化 Higgs 质量和四次耦合，首个新的六维结构来自 $u^3$：
\begin{equation*}
\Delta\widehat V_{E,H^3}^{(1)}
=\frac{\kappa^3\widehat H^3}{6}I_3(\Lambda_*),
\end{equation*}
其中
\begin{equation*}
I_3(\Lambda_*)
=\int\frac{\dd^4 k_E}{(2\pi)^4}\frac1{(k_E^2+\Lambda_*^2)^3}.
\end{equation*}
用 Schwinger 参数表示
\begin{equation*}
\frac1{A^3}=\frac12\int_0^\infty\dd\tau\,\tau^2\ee^{-\tau A}
\end{equation*}
以及四维高斯积分
\begin{equation*}
\int\frac{\dd^4 k_E}{(2\pi)^4}\ee^{-\tau k_E^2}
=\frac1{16\pi^2\tau^2},
\end{equation*}
得到
\begin{align*}
I_3(\Lambda_*)
&=\frac1{32\pi^2}\int_0^\infty\dd\tau\,\ee^{-\tau\Lambda_*^2}\\
&=\frac1{32\pi^2\Lambda_*^2}.
\end{align*}
闵可夫斯基拉格朗日密度中的势能项带负号，因此
\begin{equation*}
\Delta\widehat\Lag_{6,H}^{(1)}
=-\frac{\kappa^3}{192\pi^2\Lambda_*^2}
(\widehat\Phi^\dagger\widehat\Phi)^3.
\end{equation*}
与\eqnrefs{eq:smeft-geometric-normalization-dp13} 比较，即得到\eqnrefs{eq:heavy-scalar-ch-loop-dp32}。改回 SI Higgs 场后，系数前的整体量纲仍是 $\mathrm{J\,m^{-3}}$。
\end{derivation}

\begin{derivation}[从\eqnrefs{eq:smeft-operator-rg-dp68}到\eqnrefs{eq:smeft-wilson-evolution-dp90}]
把同一质量维数的一组重整化复合算符写成列向量 $\widehat{\mathbf O}$。正文定义
\begin{equation*}
\frac{\dd\widehat{\mathbf O}}{\dd t}
=-\gamma^{\mathcal O}(t)\widehat{\mathbf O},
\qquad t\equiv\ln\mu_{\rm DR}.
\end{equation*}
有效作用中的这一阶局域项是标量组合 $\mathbf C^T\widehat{\mathbf O}$，裸理论不能依赖辅助尺度，因此
\begin{align*}
0
&=\frac{\dd}{\dd t}(\mathbf C^T\widehat{\mathbf O})
=\left(\frac{\dd\mathbf C}{\dd t}\right)^T\widehat{\mathbf O}
+\mathbf C^T\frac{\dd\widehat{\mathbf O}}{\dd t}\\
&=\left[
\left(\frac{\dd\mathbf C}{\dd t}\right)^T
-\mathbf C^T\gamma^{\mathcal O}(t)
\right]\widehat{\mathbf O}.
\end{align*}
算符基底线性独立，所以方括号逐分量为零：
\begin{equation*}
\frac{\dd\mathbf C}{\dd t}
=(\gamma^{\mathcal O}(t))^T\mathbf C,
\end{equation*}
即正文\eqnrefs{eq:smeft-wilson-rg-dp32}。

令 $\Gamma(t)\equiv(\gamma^{\mathcal O}(t))^T$。积分方程为
\begin{equation*}
\mathbf C(t)=\mathbf C(t_0)+\int_{t_0}^{t}\dd t_1\,\Gamma(t_1)\mathbf C(t_1).
\end{equation*}
把右端的 $\mathbf C(t_1)$ 反复用同一积分方程代回，得到 Picard--Dyson 级数
\begin{align*}
\mathbf C(t)
={}&\left[I+\int_{t_0}^{t}\dd t_1\,\Gamma(t_1)
+\int_{t_0}^{t}\dd t_1\int_{t_0}^{t_1}\dd t_2\,
\Gamma(t_1)\Gamma(t_2)+\cdots\right]\mathbf C(t_0)\\
={}&\mathcal P\exp\!\left[
\int_{t_0}^{t}\dd\ell\,\Gamma(\ell)
\right]\mathbf C(t_0).
\end{align*}
第二行只是第一行嵌套积分的记号：$\mathcal P$ 保证较大的尺度参数排在矩阵乘积左侧。若 $[\Gamma(t_1),\Gamma(t_2)]=0$ 对任意 $t_1,t_2$ 都成立，嵌套积分的 $n!$ 个排序区域给出相同矩阵乘积，于是
\begin{equation*}
\mathcal P\exp\!\left[\int\dd t\,\Gamma(t)\right]
=\exp\!\left[\int\dd t\,\Gamma(t)\right].
\end{equation*}
取 $t_0=\ln\mu_{\rm M}$ 即得到正文\eqnrefs{eq:smeft-wilson-evolution-dp90}。非对角 $\gamma^{\mathcal O}$ 因而会在尺度演化中把匹配时为零的 Wilson 系数生成出来；该算符混合直接由上述矩阵微分方程的非对角演化核产生。
\end{derivation}

\begin{derivation}[从\eqnrefs{eq:weinberg-bare-renorm-dp91}到\eqnrefs{eq:weinberg-matrix-rg-dp32}]
计算两个左手轻子和两个 Higgs 的 1PI 四点函数。取 $d=4-2\epsilon$，并保留一个非零外部不变量来分开紫外与红外奇点。一圈两传播子子图都含公共泡积分
\begin{equation*}
 B(p^2)=\mu_{\rm DR}^{2\epsilon}
 \int\frac{\dd^d\ell}{(2\pi)^d}
 \frac{1}{(\ell^2+\ii0)[(\ell+p)^2+\ii0]}.
\end{equation*}
用 $1/(AB)=\int_0^1\dd x\,[xA+(1-x)B]^{-2}$，再令 $k=\ell+(1-x)p$，有
\begin{align*}
 B(p^2)
 &=\mu_{\rm DR}^{2\epsilon}\int_0^1\dd x
 \int\frac{\dd^dk}{(2\pi)^d}
 \frac{1}{[k^2-x(1-x)p^2+\ii0]^2}\\
 &=\frac{\ii}{(4\pi)^{2-\epsilon}}\Gamma(\epsilon)
 \int_0^1\dd x
 \left[\frac{\mu_{\rm DR}^2}{-x(1-x)p^2-\ii0}\right]^\epsilon.
\end{align*}
因此
\begin{equation*}
 B(p^2)\big|_{\rm UV}
 =\frac{\ii}{16\pi^2}\Delta_\epsilon,
 \qquad
 \Delta_\epsilon\equiv\frac1\epsilon-\gamma_E+\ln4\pi.
\end{equation*}
所有一圈反项都由这个极点乘上不同的味、弱同位旋与规范群系数。

味矩阵的左右位置必须由显式指标固定。带电轻子 Yukawa 项是 $\bar L_{Li}(Y_e)_{ia}\Phi e_{Ra}$，所以 Weinberg 算符左手轻子腿上的自能味矩阵是
\begin{align*}
 (\mathcal Y_L)_{ij}
 &=\sum_a(Y_e)_{ia}(Y_e^\dagger)_{aj}\\
 &=(Y_eY_e^\dagger)_{ij}.
\end{align*}
若 Yukawa 子图接在系数矩阵右侧轻子腿，味链为
\begin{align*}
 \sum_{a,\rho}C^{(5)}_{i\rho}(Y_e)_{\rho a}(Y_e^\dagger)_{aj}
 &=\sum_\rho C^{(5)}_{i\rho}(\mathcal Y_L)_{\rho j}\\
 &=(C^{(5)}\mathcal Y_L)_{ij};
\end{align*}
接在电荷共轭的左侧轻子腿时则为
\begin{align*}
 \sum_{a,\rho}(Y_e^*)_{ia}(Y_e^T)_{a\rho}C^{(5)}_{\rho j}
 &=\sum_\rho(\mathcal Y_L^T)_{i\rho}C^{(5)}_{\rho j}\\
 &=(\mathcal Y_L^TC^{(5)})_{ij}.
\end{align*}
这一步固定了 RGE 中矩阵乘法的左右次序；不能把 $Y_eY_e^\dagger$ 换成作用在右手单态空间的 $Y_e^\dagger Y_e$。

用公共因子 $(16\pi^2\epsilon)^{-1}$ 表示一圈简单极点。四点 1PI 顶角的局域反项逐类相加为
\begin{equation*}
 \delta V_5^{(1)}
 =\frac{1}{16\pi^2\epsilon}
 \left[
 \left(2\lambda_H+\frac14g'^2-\frac34g^2\right)C^{(5)}
 -\mathcal Y_L^TC^{(5)}-C^{(5)}\mathcal Y_L
 \right].
\end{equation*}
其中最后两项正是上面逐指标得到的两条 Yukawa 味链。$2\lambda_H$ 来自四 Higgs 顶角与 Weinberg 顶角的两个 Higgs 指标收缩；$g'^2/4$ 与 $-3g^2/4$ 分别来自 $U(1)_Y$ 与 $SU(2)_L$ 的规范交换及顶角子图。

还必须把四条外腿的场强反项并入同一个裸系数。左手轻子二重态和 Higgs 二重态的一圈波函数反项为
\begin{align*}
 \delta Z_L^{(1)}
 &=-\frac{1}{16\pi^2\epsilon}\frac14
 \left[g'^2+3g^2+2\mathcal Y_L\right],\\
 \delta Z_\Phi^{(1)}
 &=\frac{1}{16\pi^2\epsilon}\frac12
 \left[g'^2+3g^2-2\mathcal T_Y\right].
\end{align*}
第一式的矩阵部分仍作用在左手轻子味空间；第二式中的闭合费米子环已把味与颜色取迹，因此只剩标量 $\mathcal T_Y$。

把
\begin{equation*}
 Z_L^{-1/2}=I-\frac12\delta Z_L^{(1)}+O(g^4,Y^4),
 \qquad
 Z_\Phi^{-1}=1-\delta Z_\Phi^{(1)}+O(g^4,Y^4)
\end{equation*}
代入正文\eqnrefs{eq:weinberg-bare-renorm-dp91}。一圈 Wilson 系数反项为
\begin{equation*}
 \delta C_5^{(1)}
 =\delta V_5^{(1)}
 -\frac12(\delta Z_L^{(1)})^T C^{(5)}
 -\frac12C^{(5)}\delta Z_L^{(1)}
 -\delta Z_\Phi^{(1)}C^{(5)}.
\end{equation*}
把顶角与外腿反项按群结构逐项相加。$U(1)_Y$ 部分为
\begin{equation*}
 \frac14g'^2
 +\frac18g'^2+\frac18g'^2
 -\frac12g'^2=0,
\end{equation*}
所以一圈 Weinberg RGE 中没有 $g'^2C^{(5)}$ 项。$SU(2)_L$ 部分为
\begin{equation*}
 -\frac34g^2
 +\frac38g^2+\frac38g^2
 -\frac32g^2
 =-\frac32g^2.
\end{equation*}
带电轻子非普适部分则为
\begin{align*}
 &-\left(\mathcal Y_L^TC^{(5)}+C^{(5)}\mathcal Y_L\right)
 +\frac14\left(\mathcal Y_L^TC^{(5)}+C^{(5)}\mathcal Y_L\right)\\
 &\hspace{28mm}
 =-\frac34\left(\mathcal Y_L^TC^{(5)}+C^{(5)}\mathcal Y_L\right),
\end{align*}
而 Higgs 波函数中的 $-2\mathcal T_Y$ 经 $-\delta Z_\Phi$ 给出 $+\mathcal T_YC^{(5)}$。因此
\begin{align*}
 \delta C_5^{(1)}
 =\frac{1}{16\pi^2\epsilon}
 \Bigg\{&
 \left(2\lambda_H-\frac32g^2+\mathcal T_Y\right)C^{(5)}\\
 &-\frac34\left(\mathcal Y_L^TC^{(5)}+C^{(5)}\mathcal Y_L\right)
 \Bigg\}.
\end{align*}

简单极点转换为 beta 函数时，$A_1$ 不能当作与尺度无关的常数，因为在 $d=4-2\epsilon$ 中各无量纲耦合具有不同的工程维数。对任意一组重整化参数 $\kappa_i$，把裸量统一写成
\begin{equation*}
 \kappa_{i,B}=\mu_{\rm DR}^{\rho_i\epsilon}
 \left[\kappa_i+\frac{a_i^{(1)}(\boldsymbol\kappa)}{\epsilon}
 +\sum_{n\ge2}\frac{a_i^{(n)}(\boldsymbol\kappa)}{\epsilon^n}\right].
\end{equation*}
固定所有裸量并对 $t=\ln\mu_{\rm DR}$ 求导。$d$ 维 beta 函数分成工程维数与四维部分，
\begin{equation*}
 \frac{\dd\kappa_i}{\dd t}
 =-\rho_i\epsilon\kappa_i+\beta_i(\boldsymbol\kappa).
\end{equation*}
代回裸量的导数，在一圈阶只需保留 $a_i^{(1)}$ 中各耦合的工程维数项：
\begin{align*}
0={}&\rho_i\epsilon
 \left(\kappa_i+\frac{a_i^{(1)}}{\epsilon}\right)
 -\rho_i\epsilon\kappa_i+\beta_i
 +\frac1\epsilon\sum_j
 \left(-\rho_j\epsilon\kappa_j\right)
 \frac{\partial a_i^{(1)}}{\partial\kappa_j}
 +O((16\pi^2)^{-2})\\
={}&\beta_i+\rho_i a_i^{(1)}
 -\sum_j\rho_j\kappa_j
 \frac{\partial a_i^{(1)}}{\partial\kappa_j}
 +O((16\pi^2)^{-2}).
\end{align*}
因此 MS 类方案的一圈 simple-pole 公式是
\begin{equation*}
 \beta_i^{(1)}=
 \sum_j\rho_j\kappa_j
 \frac{\partial a_i^{(1)}}{\partial\kappa_j}
 -\rho_i a_i^{(1)}.
\end{equation*}
对 Weinberg 系数，$\rho_{C_5}=2$；当前 Higgs 势归一化下
$\rho_{\lambda_H}=2$、$\rho_g=\rho_{g'}=1$、$\rho_Y=1$。把上面的简单极点残数写成
\begin{equation*}
 a_{C_5}^{(1)}=\frac1{16\pi^2}\left[
 \left(2\lambda_H-\frac32g^2+\mathcal T_Y\right)C^{(5)}
 -\frac34\left(\mathcal Y_L^TC^{(5)}+C^{(5)}\mathcal Y_L\right)
 \right].
\end{equation*}
它对 $C^{(5)}$ 是一次齐次的。把矩阵元 $C^{(5)}_{ij}$ 视为独立坐标，Euler 定理给出
\begin{equation*}
 2\sum_{ij}C^{(5)}_{ij}
 \frac{\partial a_{C_5}^{(1)}}{\partial C^{(5)}_{ij}}
 -2a_{C_5}^{(1)}=0.
\end{equation*}
Wilson 系数自身的工程维数项因而正好抵消，剩下只需对耦合常数作齐次计数。四 Higgs 耦合项满足
\begin{equation*}
 2\lambda_H\frac{\partial}{\partial\lambda_H}
 \left(2\lambda_H C^{(5)}\right)=4\lambda_H C^{(5)},
\end{equation*}
规范项满足
\begin{equation*}
 g\frac{\partial}{\partial g}
 \left(-\frac32g^2C^{(5)}\right)=-3g^2C^{(5)}.
\end{equation*}
$\mathcal T_Y$ 与 $\mathcal Y_L$ 都是 Yukawa 矩阵的二次齐次函数；把 $Y_f$ 与 $Y_f^*$ 视为独立复变量作 Euler 计数，得到
\begin{align*}
 \sum_f\left(Y_f\!:\frac{\partial}{\partial Y_f}
 +Y_f^*\!:\frac{\partial}{\partial Y_f^*}\right)\mathcal T_Y
 &=2\mathcal T_Y,\\
 \left(Y_e\!:\frac{\partial}{\partial Y_e}
 +Y_e^*\!:\frac{\partial}{\partial Y_e^*}\right)\mathcal Y_L
 &=2\mathcal Y_L.
\end{align*}
于是 simple-pole 公式逐项给出
\begin{align*}
 16\pi^2\mu_{\rm DR}\frac{\dd C^{(5)}}{\dd\mu_{\rm DR}}
 ={}&\left(4\lambda_H-3g^2+2\mathcal T_Y\right)C^{(5)}\\
 &-\frac32\left(\mathcal Y_L^TC^{(5)}+C^{(5)}\mathcal Y_L\right),
\end{align*}
即正文\eqnrefs{eq:weinberg-matrix-rg-dp32}。右端在 $C^{(5)T}=C^{(5)}$ 时仍为复对称矩阵，所以 Majorana 系数的对称性在一圈运行下保持不变。
\end{derivation}

